\documentclass[10pt,a4paper]{amsart}
\usepackage[margin=19mm]{geometry}
\usepackage{amsmath,amssymb,mathtools}
\usepackage[scr=rsfs,cal=euler]{mathalfa}
\usepackage{xfrac}
\usepackage[unicode,hidelinks,bookmarks=false]{hyperref}
\newcommand{\ie}{i.e.\ }
\newcommand{\cf}{cf.\ }
\newcommand{\resp}{respectively\ }
\newcommand{\Subsection}{\subsection}
\providecommand{\nobreakdash}{\nobreak\hbox{-}}
\setlength{\emergencystretch}{2em}
\multlinegap0pt
\usepackage{mathrsfs}
\usepackage{xcolor}
\usepackage[shortlabels]{enumitem}
\allowdisplaybreaks

\newtheorem{definition}{Definition}
\newtheorem{theorem}{Theorem}
\newtheorem{lemma}{Lemma}
\newtheorem{fact}{Fact}
\theoremstyle{remark}
\newtheorem{assumption}{Assumption}
\newtheorem{remark}{Remark}
\newtheorem{problem}{Problem}
\newtheorem{corollary}{Corollary}
\newtheorem{example}{Example}
\newtheorem{contexample}{Counter Example}
\newtheorem{proposition}{Proposition}
\renewcommand\qedsymbol{$\blacksquare$}

\DeclareMathOperator*{\esssup}{ess\,sup}
\DeclareMathOperator*{\rank}{rank}
\DeclareMathOperator*{\spn}{span}
\DeclareMathOperator*{\im}{Im}
\DeclareMathOperator*{\proj}{proj}
\DeclareMathOperator*{\res}{res}
\DeclareMathOperator*{\kerr}{Ker}
\DeclareMathOperator*{\minimize}{minimize}

\newcommand{\Laplace}{\mathscr{L}}
\newcommand{\highlight}[1]{%
  \colorbox{yellow!50}{$\displaystyle#1$}}

\renewcommand{\thesection}{\Roman{section}}
\renewcommand{\thesubsection}{\thesection-\Alph{subsection}}
\setcounter{section}{1}
\begin{document}
\title[Complete stable inversion with necessary and sufficient conditions]{Algebraic Control: Complete Stable Inversion with Necessary and Sufficient Conditions}
\author{Burak K\"{u}rk\c{c}\"{u}, Masayoshi Tomizuka}
\date{Source: arXiv:2501.00172v4 (CC BY 4.0)}
\maketitle
\noindent\emph{Source and licence.} Algebraic Control: Complete Stable Inversion with Necessary and Sufficient Conditions. Version arXiv:2501.00172v4 (CC BY 4.0). This material is distributed under the Creative Commons Attribution 4.0 International licence, http://creativecommons.org/licenses/by/4.0/, as recorded in the arXiv OAI licence metadata for this exact version and shown on the abstract page https://arxiv.org/abs/2501.00172v4. Full rights notice: licenses/arxiv-2501.00172-CC-BY-4.0.txt.\par\smallskip
\noindent\textit{Editorial scope.} Counted unit: original Theorem 1 of the article (native environment \texttt{theorem}, source label \texttt{th:algebraic\_iff}, source0.tex lines 271--277), the necessary and sufficient conditions for exact stable inversion announced as the paper's first contribution, delivered together with its complete original author proof as the single primary proof body at source lines 278--340. The proof environment follows the statement immediately and no detached proof block is involved. No proof marker is added and no mathematics is written, repaired or re-derived. Necessary complete context is retained uncounted: source lines 131--270, that is the whole of Section II ``Preliminaries'' from its own section heading, with the subsections ``Algebraic Preliminaries'' and ``Problem Statement'', Assumption 1, Problem 1, Remark 1 and Remark 2, all of the numbered equations from (1) to (14), and the whole of the first subsection of Section III, ``Right Inverse'', up to and including the sentence that introduces the counted statement. Nothing inside the delivered ranges is omitted except the two forward cross-references described below. The publisher's own class is not present in the runtime, so the article is delivered in the AMS amsart wrapper of the pinned TeX Live runtime together with the author's own theorem declarations, the author's own math operators, the macros of the author's own preamble, and a section-numbering shim that reproduces the publisher's own numbering: \texttt{\textbackslash thesection} is set to Roman numerals and \texttt{\textbackslash thesubsection} to the form \texttt{III-C}, exactly as the version PDF prints them, with the section counter restored editorially so that the delivered section prints as Section II and the second one as Section III. The author's \texttt{rsfso} package, which the pinned runtime does not provide, is replaced by the \texttt{mathrsfs} package, which supplies the same \texttt{\textbackslash mathscr} script alphabet used for the Laplace transform symbol; no text, mathematics, numbering or source line is altered by that shim. Two forward references of Remark 2 point at subsections that lie after the counted proof and are therefore outside the delivered unit, so they are delivered as link-only adaptations under a declared \texttt{reference\_only\_adaptation} receipt: \texttt{\textbackslash ref\{sec:almostiff\}} is delivered as a link labelled with the published number \texttt{III-C} and \texttt{\textbackslash ref\{sec:practical\_modif\}} as a link labelled \texttt{III-D}, the two labels being those printed by the version PDF. The article's own bibliography is not delivered as source text: the 14 works that the delivered ranges cite are re-typeset from the author's own \texttt{.bbl} in the author's own order under the published numbers of that bibliography, so the printed citation numbers coincide with the version PDF, and the bibliography-style commands \texttt{\textbackslash newblock}, \texttt{\textbackslash bibinfo} and the \texttt{\textbackslash BIBentry} family, which only carry layout information inside those entries, are not reproduced. Printed slips kept literal, in particular the three occurrences of the opening delimiter command \texttt{\textbackslash Bigl)} in the counted statement at source lines 274 and 275, where a closing \texttt{\textbackslash Bigr)} is meant and where the printed glyph is nevertheless the closing parenthesis; the phrase ``conditions for Problem \texttt{\textbackslash ref\{pr:problem definition\}}.a. are'' at source line 272; the notation $u_z=0.y_z$ at source line 322; and the author's use of the label \texttt{pr:error}, which marks the displayed definition of the error, in the sentence ``\texttt{\textbackslash eqref\{pr:error\}} is Lyapunov-stable'' at source line 205. One declaration of the author's own preamble, \texttt{\textbackslash DeclareMathOperator*\{\textbackslash subj2\}\{subject\textbackslash to\}} at source line 57, is not carried into the wrapper: a TeX control-sequence name cannot contain a digit, so that declaration is unusable in any class and the macro could never be invoked by the name it declares; the only occurrence of that name in the whole article is at source line 404, which lies far outside the delivered unit, and every other macro of the author's preamble is reproduced verbatim. No figure, table, drawing, algorithm float or external asset occurs anywhere in the delivered ranges, so no artwork of any kind is redistributed and no bounded source comparison was needed.
\section{Preliminaries} \label{sec:Prel}
\subsection{Algebraic Preliminaries} \label{sec:algebraic background}
Let \( \mathbb{R} \) be a \textit{field} of real numbers. Consider the set of all polynomials in \( s \) with coefficients in \( \mathbb{R} \). This set forms a commutative ring over \( \mathbb{R} \) and is denoted by \( \mathbb{R}[s] \). If this ring, say \( \mathcal{R} \), has an identity element and no zero divisors, then \( \mathcal{R} \) is an \textit{integral domain}. Note that \( \mathbb{R}[s] \) is also a Euclidean Domain. 
Moreover, the set of all such rational functions in \( s \) over \( \mathbb{R} \) forms a \textit{field}, denoted by \( \mathbb{R}(s) \) \cite{forney1975minimal}, and \( \mathbb{R}[s] \subset \mathbb{R}(s) \) holds. The sets of \( n_y \times n_u \) matrices with elements in \( \mathbb{R} \), \( \mathbb{R}[s] \), and \( \mathbb{R}(s) \) are denoted by \( \mathbb{R}^{n_y \times n_u} \), \( \mathbb{R}[s]^{n_y \times n_u} \), and \( \mathbb{R}(s)^{n_y \times n_u} \) respectively.  Any $n_y \times n_u$ matrix, say $ \grave{P}(s) \in \mathbb{R}[s]^{n_y \times n_u}$, over $\mathcal{R}$ 
 can be factorized (see \textit{Invariant Factor Theorem} \cite[Theorem 2.1]{kuijper2012first}) as
 \begin{align} \label{eq:invariantfactor}
 \hspace{-0.2cm}
   \grave{P}(s)&=U_R(s)S_m(s)V_R(s) \nonumber \\   
   &=[{U}_{R1} \ {U}_{R2}]  \begin{bmatrix}
    \Lambda(s) & \textbf{0}_{r \times (n_u-r)}\\
    \textbf{0}_{(n_y-r)\times r} & \textbf{0}_{(n_y-r)\times (n_u-r)}
\end{bmatrix}{\begin{bmatrix}
    {V}_{R1} \\
    {V}_{R2}
\end{bmatrix}}
 \end{align}
where $U_R(s) \in \mathbb{R}[s]^{n_y \times n_y}$ and $V_R(s) \in \mathbb{R}[s]^{n_u \times n_u}$ are $\mathcal{R}$-Unimodular matrices, $\Lambda(s)=\text{diag}[d_1(s) \ \dots \ d_r(s)]$ with unique \textit{monic} $d_i(s) \in \mathcal{R}$ such that $d_i$ divides $d_{i+1}$. The matrix $S_m(s) \in \mathbb{R}[s]^{n_y \times n_u}$ is called the \textit{Smith Normal} form of $\grave{P}$.

Let $V$ be a vector space over the field $\mathbb{R}(s)$ with dimension $k$, consisting of $n$-tuples such that a basis of vector polynomials can always be found for \( V \) [page 22, \cite{amparan2021minimal}]. A \textit{minimal} basis of $V$ is defined as a $k \times n$ polynomial matrix $P_m$ \cite{forney1975minimal}. \textit{Adapted Forney's Theorem} \cite[Section 3(4.)]{forney1975minimal} states that if $y = xP_m$ is a polynomial $n$-tuple, then $x$ must be a polynomial $k$-tuple.

For the \textit{rank} notation, akin to \cite{antoniou2005numerical}, the rank of $P$ is defined as the maximum size of any linearly independent subset of its columns in the \textit{field} \(\mathbb{R}(s)\), denoted by \(\text{rank}_{\mathbb{R}(s)}(P)\), where \(\text{rank}_{\mathbb{R}(s)}(P) \ne \text{rank}_{\mathbb{R}}(P) \). 

Suppose $\text{rank}_{\mathbb{R}(s)}(P)=r$, where $1 \leq r \leq \min(n_u,n_y)$. {Define $J=\{j_1,\dots,j_r\} \subseteq \{1,\dots,n_u\}$ as an ordered indexed set corresponding to $P(s)$'s linearly independent columns. We then define the matrix $L$ as 
\begin{equation} \label{eq:Lind_columns_P}
{L} = [ p_{j_i} \ \dots \ p_{j_r} ]\in \mathbb{R}(s)^{n_y \times r} \, 
\end{equation}
where $p_{j_k}$ denotes the $j_k$-th column of $P$. By construction, $L$ has full column rank $r$, forming a basis for the column space of $P$ while preserving its span with a minimal set of linearly independent columns.} The image of $P$ over $\mathbb{R}(s)$ is then:
\small
\begin{align}\label{eq:spanL(P)}
\text{Im}_{\mathbb{R}(s)}(P) = \left\{ \sum_{i=1}^{r} c_i p_i : c_i \in \mathbb{R}(s), p_i \in {L} \right\} \subseteq \mathbb{R}(s)^{n_y}.
\end{align}
\normalsize
Some further notations throughout the paper are: $\Vert (\cdot)(t)\Vert$ denotes the  Euclidean norm, $\Vert (\cdot) (t) \Vert_\infty\triangleq \text{ess sup}_{t \ge 0} \Vert (\cdot)(t)\Vert$, any complex number can be expressed as $\Re(\cdot)\pm j\Im(\cdot) \in \mathbb{C}$, $\sigma$ represents singular values. Time domain square-integrable functions are denoted by $\mathcal{L}_2(-\infty, \infty)$. Its causal subset is given by $\mathcal{L}_2[0, \infty)$. For the frequency domain, including at $\infty$, $\mathcal{L}_2(j\mathbb{R})$ represents square-integrable functions on $j\mathbb{R}$, and $\mathcal{L}_\infty(j\mathbb{R})$ denotes bounded functions on $\Re(s) = 0$. All these functions in Lebesgue spaces may be either matrix-valued or scalar.
Then, $\mathcal{R}\mathcal{H}_\infty$ denotes the set of real rational $\mathcal{L}_\infty(j\mathbb{R})$ functions analytic in $\Re(s) > 0$, $(\cdot)^\dagger$ represents pseudo-inverse yielding identity matrix under multiplication, $(\cdot)^*$ denotes complex conjugates transpose, and $\langle . , . \rangle$ denotes the inner product. Among various definitions available for multivariable zeros of a Transfer Function Matrix (TFM), $P(s)$, we adopt the definition given in \cite[chapter 4.5.3]{skogestad2005multivariable} as
\begin{align} \label{eq:zerosP}
\text{Z}_P &\triangleq \{z \in \mathbb{C} : P(z)u_z=0y_z\} \\
\text{Z}_R &\triangleq  \{z \in \text{Z}_P : \Re(z) > 0\} \subseteq \text{Z}_P  \label{eq:rhpzeros}
\end{align}
where $\text{Z}_R$ defines the RHP zeros making the system non-minimum phase, $u_z,y_z$ (can be obtained via singular value decomposition (SVD) of $P(z)=U\Sigma V^*$) are normalized input and output zero directions respectively.
\subsection{Problem Statement} \label{sec:problem}
Consider the general representation of a MIMO linear time-invariant (LTI) system, $\Xi: \mathbb{U} \times \mathbb{X} \to \mathbb{Y}$, as
\begin{align} \label{eq:ss0}
\Xi(u,x_0) : \begin{cases}
\dot{x}(t)&= Ax(t)+Bu(t)\\ 
 y(t)&= C x(t), \ x_0\triangleq x(0) \ne 0
\end{cases}
\end{align} 
where $A \in \mathbb{R}^{n\times n}$, $B\in \mathbb{R}^{n\times n_u}$,  $C\in \mathbb{R}^{n_y \times n}$, $x, x_0 \in \mathbb{X} \subseteq \mathbb{R}^{n}, x(t) \in \mathcal{L}_2[0,\infty),  y \in \mathbb{Y} \subset  \mathbb{R}^{n_y}, y(t) \in \mathcal{L}_2[0,\infty), u \in \mathbb{U} \subset \mathbb{R}^{n_u}$, and $u(t) \in \mathcal{L}_2[0,\infty)$. Moreover, the TFM representation of $\Xi(u,0) \triangleq C(sI-A)^{-1}B =P(s)$ and since all real rational strictly proper transfer matrices with no poles on the
imaginary axis form $\mathcal{RL}_2(j\mathbb{R}) \subset \mathcal{L}_2(j\mathbb{R})$ \cite[page 48]{zhou1998essentials}, $ P(s) $ belongs to $ \mathcal{L}_2(j\mathbb{R})\cap \mathbb{R}(s)^{n_y \times n_u}$. The system, $\Xi$, is: P1) a \textit{non-minimum phase} s.t. $Z_R \ne \emptyset$, P2) either $n_y=n_u$ (square) or $n_y \ne n_u$ (nonsquare), and causal ($t > 0 $).
\begin{assumption}\label{ass:minimalABC}
The system given in (\ref{eq:ss0}) is minimal.
\end{assumption}

\begin{problem} \label{pr:problem definition}
{Under Assumption \ref{ass:minimalABC}, let \(\Xi\) denote the forward system given in \eqref{eq:ss0}, and let \(y_o(t)\) be an observed output trajectory. We define the \textit{right inverse} $\Xi_i^\Gamma: \mathbb{Y}\times\mathbb{X}\to\mathbb{U}$, which, given $y_o(t)$ and a prescribed inverse system's initial states, $\bar{x}_0 \in \mathbb{X}$, produces the control input
\begin{align} \label{eq:uinvpr}
u_{inv}(t) = \Xi_{i}^\Gamma(y_o(t),\bar{x}_0), \quad \Vert u_{inv}(t) \Vert_\infty < \infty.
\end{align}
Let $x_0$ be the initial states of the forward system. The error is then defined by
\begin{align} \label{pr:error}
    e(t) &\triangleq  y_o(t)-\Xi(u_{inv}(t),x_0) 
\end{align}
The inverse, $\Xi_i^\Gamma$,} is classified based on $e(t)$ as follows:
\begin{itemize}
   \item[a.]  {Stable exact inverse}: if \( e(t) = 0 \) for all \( t > 0 \).
   \item[b.]  {Stable approximate inverse}: if \( \Vert e(t) \Vert_\infty < \infty \).
\end{itemize}
Here, our goal is to construct a stabilizing inverse operator \( \Xi_i^\Gamma \)such that the error satisfies
\begin{align}\label{ESwithinevitabeerrors}
   \| e(t)\| \leq \alpha \|e(0)\| e^{-\beta t} + \varphi, \quad \forall t > 0,
\end{align}
for some $\alpha, \beta > 0$, and $\varphi \geq 0$.
\end{problem}
\begin{remark} \label{rem:similarproblems}
The terminology here (cf.\ \cite{estrada2007left}) also lets $y_o(t)$ denote a desired trajectory \cite{naderi2018inversion}.  With Assumption~\ref{ass:minimalABC} and $\|e\|_\infty<\infty$, \eqref{pr:error} is Lyapunov-stable. If exact inversion ($\alpha=\varphi=0$) is unattainable, we can employ a robust stabilizing approximate inverse—one viable choice among the many approximations acknowledged in~\cite{romagnoli2019general}—with an irreducible error~$\varphi$.
\end{remark}
\begin{remark}\label{rem:ic}  
As is typical in \textit{exact} stable inversion, we initially assume that \((A, B, C)\) in \eqref{eq:ss0} and the initial states \(x_0\) and  \(\bar{x}_0\) are known. Section~\href{https://arxiv.org/abs/2501.00172v4}{III-C} relaxes this assumption by dropping the requirement on \(\bar x_0\).
Section~\href{https://arxiv.org/abs/2501.00172v4}{III-D} goes further by allowing \(x_0\) to be unknown and by also incorporating model uncertainty, resulting in an approximate solution. For further details on initial states, see \cite{Loreto2023}.
\end{remark}

To solve Problem~\ref{pr:problem definition}.a algebraically, define 
\(y_{ic}(t) \triangleq \Xi(0, x_0)\), {$y \triangleq y_o-y_{ic}$, 
\(u_{inv}^{(0)}(t) \triangleq \Xi_i^\Gamma(0, \bar{x}_0)\), 
and \(u(t) \triangleq \Xi_i^\Gamma(y(t), 0)\) 
to represent the trajectories and control inputs for known or zero initial states. 
Then, Remark~\ref{rem:ic} lets us redefine \eqref{eq:uinvpr} as
\(u_{inv}(t) \triangleq \Xi_{i}^\Gamma(y_o(t),\bar{x}_0) - u_{inv}^{(0)}(t)\).} Applying the Laplace transform (see Paley--Wiener Th.\ \cite[p.~104]{hoffman2007banach} for existence) yields

\begin{align}
&\mathscr{L} \{ y_o(t) - {y}_{ic}(t)\} = {C}(s{I}-{A})^{-1}{B}U(s) =Y(s) \label{eq:TFM}\\
& \quad \quad \Rightarrow \ P(s)U(s)=Y(s) \quad s.t. \label{eq:algebraic equation} \\   
&\begin{bmatrix} 
P_{11}(s) & \cdots & P_{1n_u}(s) \\
\vdots & \ddots & \vdots \\
P_{n_y1}(s) & \cdots & P_{n_yn_u}(s)
\end{bmatrix}
\begin{bmatrix}
U_1(s) \\
\vdots \\
U_{n_u}(s)
\end{bmatrix}
=
\begin{bmatrix}
Y_1(s) \\
\vdots \\
Y_{n_y}(s)
\end{bmatrix}. \label{eq:algebraic equation2}
\end{align}
From this point on, similar to \cite[Theorem 8.5.2]{wolovich2012linear}, Problem \ref{pr:problem definition}.a reduces to algebraically solving the rational matrix equation where $\Vert u(t) \Vert_\infty < \infty \iff \Vert u_{inv}(t) \Vert_\infty < \infty$.
\section{Main Results} \label{sec:Main}
\subsection{{Right Inverse}} \label{sec:leftinverse}
Under Assumption \ref{ass:minimalABC}, \eqref{eq:TFM} shows that all rational function entries of $P(s)$ and $Y(s)$ in \eqref{eq:algebraic equation} share the same greatest common denominator, $\Delta=\det (sI-A) \in \mathbb{R}[s]$. To cancel out $1/\Delta$, we multiply \eqref{eq:algebraic equation} by $\Delta$, yielding 
\begin{align} \label{eq:algebraiceq_noGCD}
    \grave{P}(s)U(s)= \grave{Y}(s)
\end{align}
where $U(s) \in \mathbb{R}(s)^{n_u}$ and $\grave{Y}(s) \in \mathbb{R}(s)^{n_y}$. Then, using the \textit{invariant factor theorem} given by \eqref{eq:invariantfactor} on $\grave{P}(s)$ in \eqref{eq:algebraiceq_noGCD}, we have
\begin{align} \label{eq:leftinverse}
U(s)&=[\hat{V}_{R1} \ \hat{V}_{R2}] \begin{bmatrix}
    \Lambda^{-1}(s) & \textbf{0}_{r \times (n_y-r)}\\
    \textbf{0}_{(n_u-r)\times r} & \textbf{0}_{(n_u-r)\times (n_y-r)}
\end{bmatrix}\begin{bmatrix}
    \hat{U}_{R1} \\
    \hat{U}_{R2}
\end{bmatrix}\grave{Y}(s)\nonumber\\& \qquad +(I_{n_u}-\hat{V}_{R1}{V}_{R1})\kappa \nonumber \\
&=\Biggr( [\hat{V}_{R1} \ \hat{V}_{R2}]  \begin{bmatrix}
    \Lambda^{-1}(s) & \textbf{0}_{r \times (n_y-r)}\\
    \textbf{0}_{(n_u-r)\times r} & \textbf{0}_{(n_u-r)\times (n_y-r)}
\end{bmatrix}\begin{bmatrix}
    \hat{U}_{R1} \\
    \hat{U}_{R2}
\end{bmatrix} \nonumber\\ & \qquad +(I_{n_u}-\hat{V}_{R1}{V}_{R1})\bar{\kappa}  \Biggr)\grave{Y}(s)= \Xi_i^{\Gamma}(s) \grave{Y}(s) 
\end{align}
where $\Lambda^{-1}(s)=\text{diag}[{1}/{d_1(s)} \ \dots {1}/{d_r(s)}] \in \mathbb{R}(s)^{r \times r}$,  $\hat{U}_R(s) \in \mathbb{R}[s]^{n_y \times n_y}$ and $\hat{V}_R(s) \in \mathbb{R}[s]^{n_u \times n_u}$ are $\mathcal{R}$-Unimodular matrices such that  ${U}_R(s)\hat{U}_R(s)=I_{n_y}$ and ${V}_R(s)\hat{V}_R(s)=I_{n_u}$ respectively, and $\kappa, \bar{\kappa} \in \mathcal{R}\mathcal{H}_\infty \  (\bar{\kappa}\triangleq \kappa(\grave{Y}^T(s)\grave{Y}(s))^{-1}\grave{Y}^T(s))$ are any arbitrary vectors. 
{\begin{remark}
From an algebraic standpoint, although obtaining $U(s)$ involves multiplying $\grave{P}(s)$ by its left inverse, following \cite{moylan1977stable,estrada2007left}, we denote $\Xi_i^\Gamma$ as a “\textit{right inverse}”.
\end{remark}}
\subsection{Algebraic Necessary and Sufficient Conditions} \label{sec:algebraic_iff}
In this subsection, to find the solution(s) for Problem \ref{pr:problem definition}.a, we will show the solvability of (\ref{eq:algebraic equation}) algebraically and therefore the existence and boundedness conditions of \eqref{eq:leftinverse}. 


\begin{theorem} \label{th:algebraic_iff}
Consider $\Xi$ in \eqref{eq:ss0} with P1-P2 and $\Xi^\Gamma_i$ in \eqref{eq:leftinverse}. Then necessary and sufficient conditions for Problem \ref{pr:problem definition}.a. are
\begin{align*}
    \Bigl( e(t)=0, \ \Vert u(t) \Vert_\infty < \infty \Bigl) \iff & \Bigl( Y(s) \in \im_{\mathbb{R}(s)}(P) \Bigl) \text{ and } \\
     & \Bigl( Y(z) = 0, \ \forall z \in Z_R \Bigl).
\end{align*}
\end{theorem}

\begin{proof}
Note that $e(t)=0$ implies \eqref{eq:algebraic equation} and \eqref{eq:algebraiceq_noGCD} hold and vice versa. We can now proceed to the proof for both directions.

$\implies$Suppose $\exists U(s) \in \mathcal{R}\mathcal{H}_\infty$ satisfying $e(t)=0$.

(1:) $e(t)=0$ means \eqref{eq:leftinverse} holds. Then, pre-multiplying \eqref{eq:leftinverse} by $ \grave{P}(s) = U_R(s)S_m(s)V_R(s)$ and simplifying yields
\begin{align} \label{eq:Y invaiance}
    U_{R1}\hat{U}_{R1}\grave{Y}(s)=\grave{Y}(s)
\end{align}
which means $\grave{Y}(s)$ is invariant under the orthogonal projection onto $\im_{\mathbb{R}(s)}(P)$. Thus, we get $Y(s) \in \im_{\mathbb{R}(s)}(P)$.

(2:) Note that $(I_{n_u}-\hat{V}_{R1}{V}_{R1})$ in \eqref{eq:leftinverse} is pure polynomial and $\kappa \in \mathcal{R}\mathcal{H}_\infty$. Therefore, $(I_{n_u}-\hat{V}_{R1}{V}_{R1})\kappa$ does not contribute any unstable solution(s) to $U(s)$. So, considering only the $\hat{V}_{R1}(s)\Lambda^{-1}(s)\hat{U}_{R1}(s)\grave{Y}(s)$ is enough for stability. Let's consider two cases on $\grave{P}(s)$; C1: $\rank_{\mathbb{R}(s)}(\grave{P})=\min(n_u,n_y)$, C2: $\rank_{\mathbb{R}(s)}(\grave{P})<\min(n_u,n_y)$. 
For C2, re-writing \eqref{eq:leftinverse} and ignoring $(I_{n_u}-\hat{V}_{R1}{V}_{R1})\kappa$, we get
\begin{align}
    &U(s)=\hat{V}_{R1}\Lambda^{-1}\hat{U}_{R1}\grave{Y}(s) \implies
    \Lambda\hat{V}^\ell_{R1}U(s)=\hat{U}_{R1}\grave{Y}(s)\nonumber \\
    &\implies U_{R1}\Lambda\hat{V}^\ell_{R1} U(s)=U_{R1}\hat{U}_{R1}\grave{Y}(s) = \grave{Y}(s) \label{eq:equivalencyofnoGCD}
\end{align}
where there always exist $\hat{V}^\ell_{R1}$ satisfying $\hat{V}^\ell_{R1}\hat{V}_{R1}=I_r$. Now assume that either $\grave{P}(s)$ or $U_{R1}\Lambda\hat{V}_{R1}^\ell$ is not a minimal basis. The minimal basis \( \grave{P}_m(s) \) of an \( n_y \)-dimensional vector space of \( n_u \)-tuples over \( \mathbb{R}[s] \) for \eqref{eq:algebraiceq_noGCD} is given by \( \grave{P}_m(s) = L_T(s)\grave{P}(s) \) for condition C1, and \( \grave{P}_m(s) = L_T(s)U_{R1}\Lambda\hat{V}_{R1}^\ell \) for condition C2. Here, the \( \mathcal{R}\)-Unimodular transformation \( L_T(s) \in \mathbb{R}[s]^{n_y \times n_y} \), which yields the minimal basis, exists with constant (degree zero) determinant (see \cite[Remark 2]{forney1975minimal}). Then, left multiplying $L_T(s)$ to \eqref{eq:algebraiceq_noGCD} or \eqref{eq:equivalencyofnoGCD} results in
\begin{align} \label{eq:LTtransform_expanded}
 \grave{P}_m(s)U(s)=L_T(s)\grave{Y}(s)
\end{align}
\noindent Now define $\grave{Y}(s) \triangleq \frac{1}{\Delta_Y}{\grave{Y}_N}(s)$ such that $\grave{Y}_N(s) \in \mathbb{R}[s]^{n_y}$, and the scalar $\Delta_Y \in \mathbb{R}[s]$. So, multiplying \eqref{eq:LTtransform_expanded} by $\Delta_Y$ yields
\begin{align} \label{eq:LTtransform_expanded_purepoly}
\grave{P}_m(s)\Delta_YU(s)=L_T(s)\grave{Y}_N(s)
\end{align}
Note that $\grave{P}_m(s)$ is a minimal polynomial basis and the right-hand side of \eqref{eq:LTtransform_expanded_purepoly} is pure polynomial, so based on \textit{Adapted Forney's Theorem}, $\Delta_YU(s)$ must be polynomial. It means that $\Delta_Y$ captures all poles of $U(s)$. Moreover, non-singular transformations of $\grave{P}(s)$ preserves the invariant zeros \cite{misra1994computation} so that $y^*_z\grave{P}_m(z)=0u^*_z=0$. Left multiplying $y^*_z$ to \eqref{eq:LTtransform_expanded_purepoly} yields
\begin{align} 
\label{eq:LTtransform_expanded_purepoly_diraction}
\underbrace{y^*_z\grave{P}_m(z)}_{=0}\Delta_Y(z)U(z)=\underbrace{y^*_zL_T(z)}_{\ne 0}\grave{Y}_N(z)
\end{align}
Since the left-hand side of \eqref{eq:LTtransform_expanded_purepoly_diraction} becomes zero, $\det(L_T) = \text{constant} \ne 0$ (by polynomial-unimodularity), and $1/\Delta_Y$ is stable because $U(s) \in \mathcal{RH}_\infty$ is assumed in this direction, the only way to satisfy the equality is $\grave{Y}_N(z)=0$ implying $Y(z)=0$. 

$\impliedby$ Suppose that we have $Y(s) \in \im_{\mathbb{R}(s)}(P)$ and $Y(z)=0, \forall z \in Z_R$. For C1 \& C2, since $Y(s) \in \im_{\mathbb{R}(s)}(P)$, we have $\rank_{\mathbb{R}(s)}(P)=\rank_{\mathbb{R}(s)}([P\ : \ Y])$, which shows that we always have solution(s) for $U(s)$ in \eqref{eq:algebraiceq_noGCD} yielding $e(t)=0$ in \eqref{pr:error}, whether $U(s)$ is stable or not.

For boundedness, consider \eqref{eq:rhpzeros} and rewrite \eqref{eq:algebraiceq_noGCD} as 
\begin{align} \label{eq:algebraiceq_noGCD_Ufactored}
    \grave{P}(s)U(s) = \grave{P}(s)\frac{U_N(s)}{U_D(s)} = \grave{Y}(s) \\
    \grave{P}(s)U_N(s)=U_D(s)\grave{Y}(s) \label{eq:algebraiceq_noGCD_Ufactored2}
\end{align}
where $U_N(s) \in \mathbb{R}[s]^{n_u}$ and scalar $U_D(s) \in \mathbb{R}[s]$. Now define a new scalar $z^+(s)$ which contains all RHP zeros of $\grave{P}(s)$ as
 \begin{align*} 
    z^+(s)=\prod_{z \in Z_R \bigcap \mathbb{R}}(s - z)^{m_z}\prod_{z \in Z_R\bigcap (\mathbb{C}\setminus \mathbb{R})}(s - z)^{m_z^\ast} (s - {z}^\ast)^{m_z^\ast}
\end{align*} 
where $z^+(s) \in \mathbb{R}[s]$, $m_z$ and $m_z^\ast$ denote the multiplicities of $z$ and ${z}^*$ respectively. Thus, there exist non-zero vector $y^\ast_z \ (\text{or } u_z)$ s.t. $y^\ast_z\grave{P}(z)=0 $ (or $ \grave{P}(z)u_z=0.y_z$) $\forall z \in Z_R$. 

Now suppose $\det(U_D(\bar{z})) = 0$ for some $\bar{z}$ where $\Re(\bar{z})>0$ which means $U(s)$ is unstable. Either, suppose $(s-\bar{z}) \in z^+(s)$, then left multiplying $y^\ast_{\bar{z}}$ to \eqref{eq:algebraiceq_noGCD} and evaluating at $s=\bar{z}$ yield
\begin{align}\label{eq:contra1_th1_nec}
  \Biggl( \underbrace{(s-\bar{z})\begin{bmatrix}
\ast(s)_1 \\
\vdots \\
\ast(s)_{n_y}
\end{bmatrix}}_{= y^\ast_{\bar{z}}\grave{P}(s)}\frac{U_N(s)}{\underbrace{(s-\bar{z})U'_D(s)}_{=U_D(s)}} \Biggr) (\bar{z}) \ne 0 = y^\ast_{\bar{z}}\underbrace{\grave{Y}(\bar{z})}_{= 0}
\end{align}
or suppose $(s-\bar{z}) \notin z^+(s)$ then left multiplying $y^\ast_{\bar{z}}$ to \eqref{eq:algebraiceq_noGCD_Ufactored2} and evaluating at $s=\bar{z}$ yield
\begin{align} \label{eq:contra2_th1_nec}
  \underbrace{y^*_{\bar{z}} \grave{P}(\bar{z})}_{\ne 0}{U_N(\bar{z})} \ne 0 = {y^*_{\bar{z}}} \grave{Y}(\bar{z}) \underbrace{U_D(\bar{z})}_{= 0}
\end{align}
\eqref{eq:contra1_th1_nec} and \eqref{eq:contra2_th1_nec}
lead to a contradiction. Thus
$U(s)$ includes only LHP poles. 
The proof is now complete. 
\end{proof}

\begin{thebibliography}{99}
\bibitem[5]{moylan1977stable} P.~Moylan, ``Stable inversion of linear systems,'' \emph{IEEE Transactions on Automatic Control}, vol.~22, no.~1, pp. 74--78, 1977.
\bibitem[10]{naderi2018inversion} E.~Naderi and K.~Khorasani, ``Inversion-based output tracking and unknown input reconstruction of square discrete-time linear systems,'' \emph{Automatica}, vol.~95, pp. 44--53, 2018.
\bibitem[11]{romagnoli2019general} R.~Romagnoli and E.~Garone, ``A general framework for approximated model stable inversion,'' \emph{Automatica}, vol. 101, pp. 182--189, 2019.
\bibitem[12]{Loreto2023} M.~D. Loreto and D.~Eberard, ``Strong left inversion of linear systems and input reconstruction,'' \emph{IEEE Transactions on Automatic Control}, vol.~68, no.~6, pp. 3612--3617, 2023.
\bibitem[19]{wolovich2012linear} W.~A. Wolovich, \emph{Linear multivariable systems}.\hskip 1em plus 0.5em minus 0.4em\relax Springer Science \& Business Media, 2012, vol.~11.
\bibitem[20]{forney1975minimal} G.~D. Forney, Jr, ``Minimal bases of rational vector spaces, with applications to multivariable linear systems,'' \emph{SIAM Journal on Control}, vol.~13, no.~3, pp. 493--520, 1975.
\bibitem[21]{kuijper2012first} M.~Kuijper, \emph{First-order representations of linear systems}.\hskip 1em plus 0.5em minus 0.4em\relax Springer Science \& Business Media, 2012.
\bibitem[22]{amparan2021minimal} A.~Amparan, F.~M. Dopico, S.~Marcaida, and I.~Zaballa, ``On minimal bases and indices of rational matrices and their linearizations,'' \emph{Linear Algebra and its Applications}, vol. 623, pp. 14--67, 2021.
\bibitem[23]{antoniou2005numerical} E.~Antoniou, A.~Vardulakis, and S.~Vologiannidis, ``Numerical computation of minimal polynomial bases: A generalized resultant approach,'' \emph{Linear algebra and its applications}, vol. 405, pp. 264--278, 2005.
\bibitem[24]{skogestad2005multivariable} S.~Skogestad and I.~Postlethwaite, \emph{Multivariable feedback control: analysis and design}.\hskip 1em plus 0.5em minus 0.4em\relax john Wiley \& sons, 2005.
\bibitem[25]{zhou1998essentials} K.~Zhou and J.~C. Doyle, \emph{Essentials of robust control}.\hskip 1em plus 0.5em minus 0.4em\relax Prentice hall Upper Saddle River, NJ, 1998.
\bibitem[26]{estrada2007left} M.~B. Estrada, M.~F. Garcia, M.~Malabre, and J.~C.~M. Garc{\'\i}a, ``Left invertibility and duality for linear systems,'' \emph{Linear algebra and its applications}, vol. 425, no. 2-3, pp. 345--373, 2007.
\bibitem[27]{hoffman2007banach} K.~Hoffman, \emph{Banach spaces of analytic functions}.\hskip 1em plus 0.5em minus 0.4em\relax Courier Corporation, 2007.
\bibitem[28]{misra1994computation} P.~Misra, P.~Van~Dooren, and A.~Varga, ``Computation of structural invariants of generalized state-space systems,'' \emph{Automatica}, vol.~30, no.~12, pp. 1921--1936, 1994.
\end{thebibliography}
\end{document}
